\chapter{Resultados experimentales}
\label{ch:resultados}

Este capítulo presenta los resultados obtenidos durante la validación del prototipo de robot móvil diferencial, tanto en simulación (Gazebo) como en hardware real. Las pruebas se organizan según la progresión incremental: odometría y control diferencial, mapeo con \texttt{slam\_toolbox}, navegación autónoma con Nav2 y exploración de fronteras. Para cada etapa se describen el procedimiento, las métricas, los datos obtenidos y las observaciones derivadas, con el objetivo de cuantificar el desempeño y documentar las desviaciones entre simulación y hardware. Las pruebas de navegación y exploración, cuyos escenarios requieren un entorno de mayor extensión que el disponible, quedan planificadas según el protocolo descrito; el formato de reporte se recoge en el \ref{anx:protocolo_resultados} y se completará una vez ejecutadas en el entorno definitivo.

% ===========================================================================
\section{Descripción del entorno experimental}
\label{sec:entorno_experimental}

\subsection{Entorno de simulación}
\label{subsec:entorno_sim}

Las pruebas en simulación se realizaron en Gazebo 11 sobre Ubuntu 20.04. La configuración mantiene los mismos nombres de tópicos y estructura de parámetros que el robot real, con archivos YAML paralelos (\texttt{nav2\_params\_sim.yaml} y \texttt{mapper\_params\_online\_async\_sim.yaml}) que adaptan límites cinemáticos, frecuencias de costmap y tolerancias al entorno simulado, además de \texttt{use\_sim\_time:=true}.

\utemFigurePropia{2.Texto/Imag/gazebo_world_overview.jpg}{0.95\textwidth}
{Vista general del mundo de simulación \texttt{obstacles.world} en Gazebo.}
{fig:gazebo_world_overview}

\begin{utemMultiFigurePropia}{Secuencia de navegación en simulación (evidencia cualitativa; las métricas cuantitativas se reportan en §\ref{sec:resultados_nav2}).}{fig:sample_nav_sim_v1}
    \addutem{0.46}{2.Texto/Imag/sample_nav_sim_part1.png}{Comienzo}
    \hfill
    \addutem{0.46}{2.Texto/Imag/sample_nav_sim_part2.png}{Transición}
    \hfill
    \addutem{0.46}{2.Texto/Imag/sample_nav_sim_part3.png}{Final}
\end{utemMultiFigurePropia}

\subsection{Entorno físico de pruebas}
\label{subsec:entorno_fisico}

Las pruebas en hardware real se llevaron a cabo en un espacio interior con suelo de cerámica lisa de buena tracción. La red de comunicación se estableció mediante un \textit{travel router} con direcciones IP estáticas y \texttt{ROS\_DOMAIN\_ID=0}. La estación de trabajo (Ubuntu 20.04) ejecutó RViz2 para supervisión, mientras que toda la pila de control, mapeo y navegación se ejecutó en la Raspberry Pi 4 del robot. La extensión reducida del recinto disponible permitió validar la odometría, el control y el mapeo, pero limitó los escenarios de navegación y exploración, que requieren trayectorias y áreas mayores.

\subsection{Instrumentación y métricas}
\label{subsec:metricas}

Para la evaluación cuantitativa se definieron las métricas de la Tabla~\ref{tab:metricas_experimentales}.

\begin{table}[H]
\centering
\begin{threeparttable}
\caption{Métricas de evaluación utilizadas.}
\label{tab:metricas_experimentales}
\small
\begin{tabular}{p{3.8cm} p{1.5cm} p{5cm} p{2.2cm}}
\toprule
\textbf{Métrica} & \textbf{Unidad} & \textbf{Herramienta} & \textbf{Aplicación} \\
\midrule
Error de posición en línea recta & m & Medición física + odometría & Odometría \\
Deriva angular en giro de 360° & grados & Medición visual + \texttt{/odom} & Odometría \\
Tiempo de respuesta al escalón\tnote{a} & s & Modelo FODT identificado & Control PID \\
Sobreimpulso & \% & Análisis de la respuesta & Control PID \\
Error en estado estacionario & m/s & \texttt{/motor\_vels} & Control PID \\
Tiempo de construcción del mapa & s & Cronometraje + \texttt{ros2 bag} & SLAM \\
Consistencia del mapa & --- & Inspección visual en RViz & SLAM \\
Cierre de lazo & --- & Alineación de paredes en RViz & SLAM \\
Tasa de éxito de navegación & \% & Metas alcanzadas / enviadas & Nav2 \\
Tiempo de navegación por meta & s & Cronometraje envío--arribo & Nav2 \\
Cobertura de exploración & \% & Área mapeada / área total & Exploración \\
\bottomrule
\end{tabular}
\begin{tablenotes}
\footnotesize
\item[a] Calculado mediante simulación sobre el modelo FODT identificado, no por medición directa sobre el hardware.
\end{tablenotes}
\end{threeparttable}
\end{table}

\paragraph{Definición formal de las métricas del controlador PID.}
Los indicadores se calculan sobre la respuesta $y[k]$ a un escalón $r$:

\begin{align}
M_p &= \max\left(0,\;\frac{\max_k y[k] - r}{r} \times 100\right) \quad [\%]
\label{eq:overshoot} \\[2pt]
t_r &= t_{90} - t_{10},\quad t_{\alpha} = \min\{t[k] : y[k] \geq \alpha r\},\; \alpha\in\{0.1,0.9\}
\label{eq:rise_time} \\[2pt]
t_s &= t[k^{*}+1],\quad k^{*} = \max\{k : |y[k] - r| > 0.02\,r\}
\label{eq:settling_time} \\[2pt]
e_{ss} &= \left|r - \frac{1}{N - N_{ss}}\sum_{k=N_{ss}}^{N-1} y[k]\right|,\quad N_{ss} = \lfloor 0.9N\rfloor
\label{eq:ess}
\end{align}

\noindent
correspondientes a sobreimpulso máximo ($M_p$), tiempo de subida ($t_r$, 10\%--90\%), tiempo de asentamiento ($t_s$, banda 2\%) y error en estado estacionario ($e_{ss}$). Las métricas comparativas PSO vs. Lambda se obtienen con \texttt{deadband}$\,=0$ por consistencia metodológica. Los datos de odometría, velocidad y transformaciones se registraron con \texttt{ros2 bag record}.

% ===========================================================================
\section{Validación de odometría y control diferencial}
\label{sec:resultados_odometria}

\subsection{Prueba de trayectoria rectilínea}
\label{subsec:prueba_recta}

\textbf{Procedimiento:} se envió una velocidad lineal constante de $0.13$ m/s durante el tiempo correspondiente a $2.00$ m teóricos, con cinco repeticiones; se midió la distancia real con cinta métrica y se comparó con \texttt{/odom}. El error porcentual se define como:

\begin{equation}
e_{\text{odom}} = \frac{|d_{\text{teórica}} - d_{\text{medida}}|}{d_{\text{teórica}}} \times 100 \quad [\%]
\label{eq:error_odom}
\end{equation}

\noindent
En hardware real, $d_{\text{medida}}$ es la distancia física medida con cinta, por lo que el indicador cuantifica el error de odometría propiamente tal (discrepancia entre encoders y realidad). En simulación, la odometría del plugin de Gazebo coincide con la verdad de terreno (\textit{ground truth}), de modo que el indicador cuantifica el error de seguimiento de consigna (diferencia entre distancia comandada y ejecutada por el controlador), no un error de odometría en sentido estricto.

\begin{table}[H]
\centering
\begin{threeparttable}
\caption{Error de seguimiento de consigna en trayectoria rectilínea -- simulación (Gazebo).}
\label{tab:odom_recta_sim}
\begin{tabular}{ccccc}
\toprule
\textbf{Prueba} & \textbf{Teórica (m)} & \textbf{Odom. (m)} & \textbf{Real GT (m)}\tnote{a} & \textbf{Error (\%)} \\
\midrule
1 & 2.00 & 1.866 & 1.866 & 6.69 \\
2 & 2.00 & 1.954 & 1.954 & 2.28 \\
3 & 2.00 & 1.946 & 1.946 & 2.72 \\
4 & 2.00 & 1.923 & 1.923 & 3.83 \\
5 & 2.00 & 1.972 & 1.972 & 1.40 \\
\midrule
\multicolumn{4}{r}{\textbf{Error medio}} & \textbf{3.39} \\
\bottomrule
\end{tabular}
\begin{tablenotes}
\footnotesize
\item[a] En Gazebo la odometría coincide con el \textit{ground truth}; se reporta una sola columna.
\end{tablenotes}
\end{threeparttable}
\end{table}

El error en simulación no refleja imprecisión de la odometría —ideal en Gazebo— sino la diferencia entre la distancia comandada y la efectivamente recorrida, atribuible a la dinámica del controlador (aceleración, frenado y discretización). La primera prueba presenta mayor error (6.69\%), atribuible a transitorios iniciales antes de alcanzar la referencia; el error medio sin ella es 2.56\%.

\utemFigurePropia{2.Texto/Imag/prueba_recta.png}{0.90\textwidth}
{Trayectoria rectilínea en Gazebo.}
{fig:prueba_recta}

\begin{table}[H]
\centering
\begin{threeparttable}
\caption{Error de odometría en trayectoria rectilínea -- hardware real.}
\label{tab:odom_recta_real}
\begin{tabular}{ccccc}
\toprule
\textbf{Prueba} & \textbf{Teórica (m)} & \textbf{Odom. (m)} & \textbf{Real (m)} & \textbf{Error (\%)} \\
\midrule
1 & 2.00 & 1.920 & 1.846 & 4.01 \\
2 & 2.00 & 1.922 & 1.868 & 2.88 \\
3 & 2.00 & 1.923 & 1.851 & 3.87 \\
4 & 2.00 & 1.926 & 1.856 & 3.78 \\
5 & 2.00 & 1.928 & 1.861 & 3.61 \\
\midrule
\multicolumn{4}{r}{\textbf{Error medio}} & \textbf{3.63} \\
\bottomrule
\end{tabular}
\begin{tablenotes}
\footnotesize
\item En hardware, $d_{\text{medida}}$ es la distancia física con cinta; el error cuantifica la discrepancia encoders--realidad (Ec.~\ref{eq:error_odom}).
\end{tablenotes}
\end{threeparttable}
\end{table}

La distancia reportada por \texttt{/odom} es sistemáticamente inferior a la teórica, lo que sugiere que el controlador no alcanza plenamente la velocidad de referencia por el límite de saturación ($0.13$ m/s).

\subsection{Prueba de giro completo (360°)}
\label{subsec:prueba_giro}

\textbf{Procedimiento:} se aplicó una velocidad angular pura de $0.4$ rad/s (tiempo teórico $2\pi/0.4 \approx 15.71$ s), marcando la orientación inicial y midiendo el ángulo residual; cinco repeticiones.

\begin{table}[H]
\centering
\begin{threeparttable}
\caption{Error de seguimiento en giro completo -- simulación.}
\label{tab:odom_giro_sim}
\begin{tabular}{ccccc}
\toprule
\textbf{Prueba} & \textbf{Teórico (\textdegree)} & \textbf{Odom. (\textdegree)} & \textbf{Real GT (\textdegree)} & \textbf{Error (\%)}\tnote{a} \\
\midrule
1 & 360.0 & 335.86 & 335.86 & 6.71 \\
2 & 360.0 & 347.40 & 347.40 & 3.50 \\
3 & 360.0 & 346.62 & 346.62 & 3.72 \\
4 & 360.0 & 347.40 & 347.40 & 3.50 \\
5 & 360.0 & 337.41 & 337.41 & 6.28 \\
\midrule
\multicolumn{4}{r}{\textbf{Error medio}} & \textbf{4.74} \\
\bottomrule
\end{tabular}
\begin{tablenotes}
\footnotesize
\item[a] Error $= |360^\circ - \text{GT}| / 360^\circ \times 100$. El ángulo teórico corresponde al tiempo de comando, no a la lectura de \texttt{/odom}.
\end{tablenotes}
\end{threeparttable}
\end{table}

\utemFigurePropia{2.Texto/Imag/prueba_giro.png}{0.90\textwidth}
{Simulación del giro de 360°.}
{fig:prueba_giro}

\begin{table}[H]
\centering
\begin{threeparttable}
\caption{Deriva angular tras giro completo de 360° -- hardware real.}
\label{tab:odom_giro}
\begin{tabular}{cccc}
\toprule
\textbf{Prueba} & \textbf{Reportado (°)} & \textbf{Real (°)} & \textbf{Deriva (°)} \\
\midrule
1 & 360.0 & 355.7 & $-4.3$ \\
2 & 360.0 & 340.0 & $-20.0$ \\
3 & 360.0 & 380.0 & $+20.0$ \\
4 & 360.0 & 370.0 & $+10.0$ \\
5 & 360.0 & 356.0 & $-4.0$ \\
\midrule
\multicolumn{3}{r}{\textbf{Deriva media ($|\text{Deriva}|$)}} & \textbf{11.7} \\
\bottomrule
\end{tabular}
\begin{tablenotes}
\footnotesize
\item Deriva $=$ ángulo real medido $-$ ángulo reportado por \texttt{/odom} ($360^\circ$).
\end{tablenotes}
\end{threeparttable}
\end{table}

La alta variabilidad apunta a factores no controlados (superficie, zona muerta del L298N). Excluyendo las dos repeticiones de mayor deriva (pruebas 2 y 3, $\pm 20^\circ$), la deriva media de las restantes es $6.1^\circ$. En ambos casos, la magnitud refuerza la necesidad de una IMU para corregir la orientación, consistente con la decisión de no activar AMCL (Capítulo~\ref{chap:app_control_nav}).

\subsection{Cadena de transformaciones \texttt{tf}}
\label{subsec:verificacion_tf}

Se verificó que el árbol de transformaciones se publica de forma continua (\texttt{map}$\rightarrow$\texttt{odom} por \texttt{slam\_toolbox}; \texttt{odom}$\rightarrow$\texttt{base\_link} por \texttt{diff\_drive\_controller}; \texttt{base\_link}$\rightarrow$\texttt{laser\_frame} estática), y que \texttt{/odom} se publica a $\approx 30$ Hz, limitado por el \texttt{update\_rate} del \emph{controller manager} (Figura~\ref{fig:tf_tree_real}).

\utemFigurePropia{2.Texto/Imag/tf_tree_real.pdf}{0.88\textwidth}
{Árbol de transformaciones del robot real.}
{fig:tf_tree_real}

% ===========================================================================
\section{Evaluación del controlador PID sintonizado}
\label{sec:resultados_pid}

\subsection{Respuesta al escalón con ganancias finales}
\label{subsec:respuesta_escalon_pso}

Se aplicó un escalón de velocidad de $0.13$ m/s con las ganancias definitivas del firmware y se registró la respuesta experimental (Figura~\ref{fig:respuesta_escalon_final}). Los indicadores de desempeño se calculan sobre el modelo FODT identificado a partir de esa captura, utilizando dichas ganancias; la Tabla~\ref{tab:fodt_params_pso} presenta el modelo identificado y la Tabla~\ref{tab:pid_performance} los indicadores resultantes.

\utemFigurePropia{2.Texto/Imag/respuesta_escalon_final.png}{0.88\textwidth}
{Respuesta experimental al escalón con ganancias finales.}
{fig:respuesta_escalon_final}

\begin{table}[H]
\centering
\begin{threeparttable}
\caption{Parámetros FODT identificados mediante PSO sobre la respuesta en lazo abierto.}
\label{tab:fodt_params_pso}
\begin{tabular}{lccc}
\toprule
\textbf{Rueda} & $K_m$ (m/s/PWM) & $\tau$ (s) & $L_d$ (s)\tnote{a} \\
\midrule
Izquierda & 0.000913 & 0.3291 & 0.0215 / 0.0001 \\
Derecha   & 0.000890 & 0.3485 & 0.0001 / 0.0845 \\
\bottomrule
\end{tabular}
\begin{tablenotes}
\footnotesize
\item[a] El PSO es estocástico: $K_m$ y $\tau$ son consistentes entre corridas, mientras que el retardo $L_d$ varía. Se reportan los valores de las dos corridas empleadas, respectivamente, para la sintonía PSO y para la comparación con Lambda.
\end{tablenotes}
\end{threeparttable}
\end{table}

\begin{table}[H]
\centering
\begin{threeparttable}
\caption{Indicadores de desempeño PID con las ganancias finales (modelo FODT, \texttt{deadband}$\,=0$).}
\label{tab:pid_performance}
\begin{tabular}{lcc}
\toprule
\textbf{Indicador} & \textbf{Rueda izquierda} & \textbf{Rueda derecha} \\
\midrule
Tiempo de subida (10\%--90\%) & 0.600~s & 0.600~s \\
Tiempo de establecimiento (2\%) & N/A\tnote{a} & 3.100~s \\
Sobreimpulso máximo & 15.63\% & 10.66\% \\
Error estacionario & 0.00488~m/s & 0.00192~m/s \\
\bottomrule
\end{tabular}
\begin{tablenotes}
\footnotesize
\item[a] La rueda izquierda no alcanza la banda del 2\% dentro de 10~s por oscilación residual.
\end{tablenotes}
\end{threeparttable}
\end{table}

\subsection{Comparación con sintonía Lambda}
\label{subsec:comparacion_lambda}

Se calcularon ganancias Lambda ($\lambda = 3\tau$) sobre el mismo modelo FODT. La Tabla~\ref{tab:comparacion_metodos} compara ambos enfoques.

\begin{table}[H]
\centering
\begin{threeparttable}
\caption{Comparativa de ganancias PID\tnote{a}: firmware (PSO + refinamiento) vs. Lambda.}
\label{tab:comparacion_metodos}
\footnotesize
\begin{tabular}{l ccc ccc}
\toprule
\textbf{Método} & $k_p$ & $k_i$ & $k_d$ & \textbf{Sobreimp. (\%)} & $t_r$ (s) & $e_{ss}$ (m/s) \\
\midrule
Firmware (Izq.) & 0.7017 & 0.09867 & 0.8902 & 15.63 & 0.600 & 0.00488 \\
Lambda (Izq.)   & 2.7193 & 7.50238 & 0.0823 & 46.42 & 0.200 & 0.00015 \\
Firmware (Der.) & 0.7201 & 0.00132 & 0.8793 & 10.66 & 0.600 & 0.00192 \\
Lambda (Der.)   & 2.8241 & 7.22823 & 0.1064 & 44.15 & 0.200 & 0.00005 \\
\bottomrule
\end{tabular}
\begin{tablenotes}
\scriptsize
\item[a] Ganancias efectivas en escala continua ($k_p=K_p/K_o$, etc.) para comparar con Lambda: los términos proporcional y derivativo se escalan por $f_\lambda \approx 22.80$, y el integral incorpora además el período de muestreo del lazo.
\end{tablenotes}
\end{threeparttable}
\end{table}

Lambda es más rápido (menor $t_r$ y $e_{ss}$), pero con un sobreimpulso inaceptable ($>40\%$). El criterio PSO+experimental prioriza el bajo sobreimpulso (10.66\%--15.63\%) a costa de un mayor tiempo de subida, comportamiento preferible para un robot que opera cerca de obstáculos.

\utemFigurePropia{2.Texto/Imag/comparacion_pso_vs_lambda.png}{0.92\textwidth}
{Comparación simulada PSO vs. Lambda sobre el modelo FODT.}
{fig:comparacion_pso_vs_lambda}

\subsection{Trayectoria recta a velocidad constante}
\label{subsec:pid_trayectoria}

Durante un recorrido recto de 3 m a $0.13$ m/s, ambas ruedas siguieron la referencia con oscilaciones de baja amplitud y mantuvieron una trayectoria aproximadamente rectilínea, lo que confirma que las ganancias asimétricas compensan las diferencias mecánicas entre motores.

% ===========================================================================
\section{Resultados de mapeo con \texttt{slam\_toolbox}}
\label{sec:resultados_slam}

\subsection{Mapeo en simulación}
\label{subsec:mapa_sim}

Se teleoperó el robot en \texttt{obstacles.world} mientras \texttt{slam\_toolbox} construía el mapa (Tabla~\ref{tab:slam_sim_indicadores}). El mapa resultante presenta paredes nítidas y obstáculos bien representados.

\begin{table}[H]
\centering
\begin{threeparttable}
\caption{Indicadores del mapeo en simulación.}
\label{tab:slam_sim_indicadores}
\begin{tabular}{ll}
\toprule
\textbf{Indicador} & \textbf{Valor} \\
\midrule
Tiempo total de mapeo & 10.1 min \\
Distancia recorrida (odometría) & 23.32 m \\
Número de nodos del grafo & 45 \\
Resolución del mapa & 0.05 m/celda \\
Cobertura estimada & 95.6\% \\
Área libre mapeada & 42.53 m² \\
Cierre de lazo & Consistencia global correcta\tnote{a} \\
Alineación de paredes & Bordes definidos \\
\bottomrule
\end{tabular}
\begin{tablenotes}
\footnotesize
\item[a] El cierre de lazo no generó un evento explícito, pero la consistencia global fue buena.
\end{tablenotes}
\end{threeparttable}
\end{table}

\subsection{Mapeo en hardware real}
\label{subsec:mapa_real}

El procedimiento se repitió en el entorno físico, obteniéndose mapas utilizables. La calidad depende críticamente de la velocidad de desplazamiento: las mejores capturas se lograron entre 0.06 y 0.09 m/s, rango que permite al \gls{LIDAR} completar un barrido entre actualizaciones del grafo. Los indicadores cuantitativos (tiempo, distancia, nodos del grafo y cierre de lazo) quedan pendientes de medición sistemática; el formato de reporte, idéntico al de la Tabla~\ref{tab:slam_sim_indicadores}, se recoge en el \ref{anx:protocolo_resultados}.

\subsection{Comparación cualitativa de mapas}
\label{subsec:comparacion_mapas}

\begin{table}[H]
\centering
\caption{Diferencias cualitativas observadas entre mapas de simulación y hardware real.}
\label{tab:diferencias_mapas}
\begin{tabular}{lll}
\toprule
\textbf{Aspecto} & \textbf{Simulación} & \textbf{Hardware real} \\
\midrule
Nitidez de paredes & Alta (bordes definidos) & Media \\
Ruido en espacio libre & Nulo & Bajo--moderado \\
Consistencia en cierre & Alta & Alta/media \\
Obstáculos pequeños & Detectados & Parcialmente detectados \\
Zonas de sombra láser & Mínimas & Presentes \\
\bottomrule
\end{tabular}
\end{table}

% ===========================================================================
\section{Resultados de navegación autónoma con Nav2}
\label{sec:resultados_nav2}

\subsection{Protocolo de pruebas}
\label{subsec:protocolo_nav}

Se definieron tres escenarios (cinco repeticiones cada uno): meta directa sin obstáculos ($\approx 3$ m), meta con obstáculo intermedio (caja) y meta en pasillo/puerta (espacio estrecho). La meta se envía desde RViz y se considera éxito si el robot la alcanza dentro de la tolerancia \texttt{xy\_goal\_tolerance} (0.25 m en simulación, 0.22 m en hardware). Estos escenarios exceden la extensión del recinto físico disponible, por lo que su ejecución queda planificada para un entorno de mayor superficie; el formato de reporte (tasa de éxito, tiempo medio y desviación por escenario, en simulación y hardware) se recoge en el \ref{anx:protocolo_resultados}.

\subsection{Tipos de fallo y estrategias de recuperación}
\label{subsec:analisis_fallos}

La Tabla~\ref{tab:fallos_nav} clasifica los modos de fallo previstos y la acción de recuperación del árbol de comportamiento que los atiende; la tasa de recuperación observada se reportará junto con los resultados de navegación.

\begin{table}[H]
\centering
\begin{threeparttable}
\caption{Tipos de fallo y estrategias de recuperación.}
\label{tab:fallos_nav}
\small
\begin{tabular}{p{3.5cm} p{4cm} p{4.5cm}}
\toprule
\textbf{Tipo de fallo} & \textbf{Causa probable} & \textbf{Recuperación} \\
\midrule
Costmap bloqueado & Obstáculo no mapeado & \texttt{ClearCostmap}\tnote{a} \\
Ruta inválida & Espacio muy estrecho & Replaneación (4~Hz) \\
Robot atascado & Deslizamiento en giro & \texttt{Spin} + \texttt{BackUp} \\
Timeout de navegación & Oscilación local & Abortar y reintentar \\
\bottomrule
\end{tabular}
\begin{tablenotes}
\footnotesize
\item[a] Limpieza en radio 1.0~m (local) y 1.5~m (global).
\end{tablenotes}
\end{threeparttable}
\end{table}

% ===========================================================================
\section{Resultados de exploración autónoma de fronteras}
\label{sec:resultados_exploracion}

La exploración se evalúa colocando el robot en una esquina y lanzando \texttt{explore\_server}, registrando el número de metas generadas y alcanzadas, la cobertura y las intervenciones manuales. Al igual que la navegación, requiere un área mayor que la disponible, por lo que su ejecución queda planificada; el formato de reporte se recoge en el \ref{anx:protocolo_resultados}.

% ===========================================================================
\section{Coherencia simulación -- hardware real}
\label{sec:comparacion_sim_real}

\begin{table}[H]
\centering
\begin{threeparttable}
\caption{Comparación integral de desempeño: simulación vs. hardware real.}
\label{tab:comparacion_integral}
\small
\begin{tabular}{p{3.2cm} p{3.8cm} p{3cm} p{3cm}}
\toprule
\textbf{Subsistema} & \textbf{Métrica} & \textbf{Simulación} & \textbf{Hardware real} \\
\midrule
Odometría (recta 2 m)\tnote{b} & Error medio (\%) & 3.39 & 3.63 \\
Odometría (giro 360°)\tnote{c} & Desviación angular & $\approx 17.1^\circ$ (seguim.) & $11.7^\circ$ (deriva) \\
Control PID\tnote{a} & Sobreimpulso (\%) & N/A & 15.63 / 10.66 \\
Control PID\tnote{a} & Error estacionario (m/s) & N/A & 0.00488 / 0.00192 \\
Mapeo SLAM & Nitidez de paredes & Alta & Media--alta \\
Mapeo SLAM & Cierre de lazo & Consistente & Consistente \\
Navegación Nav2 & Tasa de éxito (\%) & Pendiente & Pendiente \\
Navegación Nav2 & Tiempo medio a meta (s) & Pendiente & Pendiente \\
Exploración & Cobertura (\%) & Pendiente & Pendiente \\
\bottomrule
\end{tabular}
\begin{tablenotes}
\footnotesize
\item[a] En simulación, Gazebo gestiona la dinámica de los motores mediante su motor de física, por lo que el PID del firmware no interviene y no es evaluable. Los indicadores de la columna ``hardware real'' provienen del modelo FODT identificado experimentalmente con las ganancias que operan en el robot (Sección~\ref{sec:resultados_pid}).
\item[b] No estrictamente comparable: en Gazebo el error refleja el seguimiento de consigna (odometría $=$ \textit{ground truth}); en hardware refleja la discrepancia entre encoders y la medición física.
\item[c] Tampoco comparable directamente: en simulación cuantifica el seguimiento de consigna (el robot recorre $\approx 343^\circ$ de los $360^\circ$ comandados por dinámica y discretización); en hardware cuantifica la deriva de la odometría respecto de la orientación física. Se reportan en sus unidades nativas y no deben restarse.
\end{tablenotes}
\end{threeparttable}
\end{table}

Respecto de los criterios de paridad definidos al construir el gemelo digital (Capítulo~\ref{ch:robot_digital}), la coherencia cinemática se verifica: la separación de ruedas, el radio y la base de conversión de odometría son idénticos en ambos entornos, y el error de odometría se mantiene en una banda equivalente (3.39\% en simulación, 3.63\% en hardware). El cierre cuantitativo de los criterios asociados a navegación se completará con los resultados pendientes.

\subsection{Adaptaciones en la transferencia simulación -- hardware}
\label{subsec:adaptaciones_transfer}

La transferencia al hardware requirió fijar \texttt{use\_sim\_time:=false}, activar las reglas \texttt{udev} de nombres persistentes, ajustar los límites cinemáticos y las frecuencias del controlador, y habilitar el retroceso. El detalle completo de las diferencias de configuración entre entornos se documenta en los Capítulos~\ref{chap:robot_real} y~\ref{chap:app_control_nav}.

% ===========================================================================
\section{Viabilidad en la Raspberry Pi 4}
\label{sec:recursos_computacionales}

Durante las sesiones de SLAM y navegación superiores a 10~min de operación continua no se observaron caídas de tensión ni sobrecalentamiento en el riel de 5~V, comportamiento consistente con el margen del convertidor frente al presupuesto de potencia (\ref{anx:hardware_ros2}); la etapa de potencia opera con reserva suficiente para la carga del prototipo y la incorporación futura de sensores. La caracterización cuantitativa del uso de cómputo y de la autonomía energética queda pendiente de medición sistemática.

% ===========================================================================
\section*{Conclusión del capítulo}
\label{sec:conclusion_resultados}

La odometría basada en encoders de 1980 CPR ofrece precisión adecuada para trayectorias cortas (error medio 3.63\% en 2 m en hardware real), pero la deriva angular media de $11.7^\circ$ por giro completo se propaga al mapa, confirmando la necesidad de una IMU. La sintonización PSO+experimental produce bajo sobreimpulso (10.66\%--15.63\%) a costa de un mayor tiempo de subida; el \emph{ripple} residual ($\approx 0.005$ m/s) proviene de la cuantización del encoder y no afecta la navegación, y la zona muerta asimétrica suprime los ciclos límite. El mapeo SLAM produce mapas utilizables en ambos entornos; en hardware la calidad depende de mantener velocidades bajas (0.06--0.09 m/s) para que el \gls{LIDAR} complete un barrido entre actualizaciones del grafo. La etapa de potencia se mostró estable durante toda la experimentación, sin síntomas de sobrecarga ni estrangulamiento térmico.

La validación de la navegación autónoma y la exploración de fronteras quedó limitada por la extensión del recinto físico disponible, que no admite los escenarios definidos (meta a $\approx 3$ m, pasillo). El protocolo, las métricas y la infraestructura de pruebas están definidos (\ref{anx:protocolo_resultados}), de modo que estas evaluaciones pueden completarse en un entorno de mayor superficie para cerrar cuantitativamente la validación del sistema integrado.